\section{Introduction}


Blockchain consensus mechanisms face a fundamental challenge: how to fairly allocate block production rights and rewards. Proof of Work (PoW) solves this by requiring computational effort, but the computation serves no purpose beyond securing the chain---an estimated 150 TWh annually for Bitcoin alone. Proof of Stake (PoS) reduces energy waste but creates plutocratic dynamics where the wealthy accumulate more wealth.

The rise of AI compute networks presents an opportunity to align consensus with useful work. AI training and inference require substantial compute resources; if this compute could simultaneously secure a blockchain, we would achieve both useful work and decentralized consensus.

\subsection{Challenges}

Designing \PoAI{} requires solving several challenges:

\begin{enumerate}
    \item \textbf{Verification}: How can validators confirm that claimed AI work was actually performed?

    \item \textbf{Quality Assessment}: Not all AI compute is equally valuable---training a useful model differs from generating noise.

    \item \textbf{Sybil Resistance}: Without proof of identity, adversaries could claim credit for others' work.

    \item \textbf{Efficiency}: Verification must be cheaper than the original computation.

    \item \textbf{Incentive Compatibility}: Rational nodes must be incentivized to perform high-quality work.
\end{enumerate}

\subsection{Contributions}

This paper makes the following contributions:

\begin{enumerate}
    \item \textbf{\PoAI{} Protocol}: A complete consensus mechanism combining \TEE{} attestation, quality scoring, and BLS aggregation (Section~\ref{sec:protocol})

    \item \textbf{Compute Verification}: Novel techniques for verifying AI workloads without re-executing them (Section~\ref{sec:verification})

    \item \textbf{Quality Scoring}: Multi-party evaluation protocol for assessing AI output value (Section~\ref{sec:quality})

    \item \textbf{Security Analysis}: Proofs of Byzantine fault tolerance and incentive compatibility (Section~\ref{sec:security})

    \item \textbf{Evaluation}: Experimental results demonstrating scalability and efficiency (Section~\ref{sec:evaluation})

    \item \textbf{Execution-level proofs (software-only)}: We close the gap that
    signature-plus-plurality leaves open---paying for a guess rather than for
    computation (Section~\ref{sec:guess-to-mint})---with one primitive yielding four
    proofs (Section~\ref{sec:four-proofs}), a software-only Freivalds verifier with a
    soundness theorem (Section~\ref{sec:freivalds}), a domain-separated task binding
    with anti-replay and anti-substitution lemmas (Section~\ref{sec:binding}), the
    Bitcoin shared-difficulty correspondence (Section~\ref{sec:difficulty}), a
    single-gate tier matrix (Section~\ref{sec:tiers}), and a user-sovereignty trust
    spectrum (Section~\ref{sec:ownership}).

    \item \textbf{Full modern-LLM coverage}: We show that ``support every modern LLM''
    is achieved by \emph{tiering} verification to the arithmetic---never by a
    floating-point Freivalds, which is unsound (a tolerance corridor) and unusable
    (cross-hardware false-reject) at once---and that floating-point GEMMs become
    bit-exact by pinning the reduction order (Section~\ref{sec:fp-coverage}). We map the
    \Zoo{} model zoo (dense, mixture-of-experts, linear-recurrent/state-space,
    multimodal vision, audio, and diffusion) to proof obligations and prove that beyond
    base matmul checking exactly \emph{two} new core techniques suffice---a routing proof
    and a sequential-chaining proof (Section~\ref{sec:families}). We report the build
    status honestly: the Freivalds verifier core and the two determinism prerequisites
    are implemented and unit-tested in the engine, while the commitment layer and the
    two extensions are designed (Section~\ref{sec:impl-status}).
\end{enumerate}

